\documentclass[a4paper]{article}
\usepackage{amsfonts}
\usepackage{amsmath}
\usepackage{amssymb}
\usepackage{graphicx}     

\newcommand{\Bgtan}{\operatorname{Bgtan}}

\begin{document}
  
\noindent \large Auteur: Siebe Mees \\
\noindent \large Studierichting: Informatica\\
\noindent \large Oefeningenreeks 2B nr. 6.7\\

\medskip

\normalsize


$4\Bgtan\dfrac{1}{5} - \Bgtan\dfrac{1}{239} = \dfrac{\pi}{4}$\\

$= 2\Bgtan\dfrac{\dfrac{1}{5}+\dfrac{1}{5}}{1-\dfrac{1}{5}\cdot\dfrac{1}{5}} - \Bgtan\dfrac{1}{239}$\\

$= 2\Bgtan\dfrac{10}{24} - \Bgtan\dfrac{1}{239}$\\

$= \Bgtan\dfrac{\dfrac{10}{24}+\dfrac{10}{24}}{1-\dfrac{10}{24}\cdot\dfrac{10}{24}} - \Bgtan\dfrac{1}{239}$\\

$= \Bgtan\dfrac{120}{119} - \Bgtan\dfrac{1}{239}$\\

$= \Bgtan\dfrac{\dfrac{120}{119}-\dfrac{1}{239}}{1+\dfrac{120}{119}\cdot\dfrac{1}{239}} $\\

$= \Bgtan\dfrac{\dfrac{28680}{28441}-\dfrac{119}{28441}}{\dfrac{28441}{28441}+\dfrac{120}{28441}} $\\

$= \Bgtan1$\\

\begin{thebibliography}{999}
%\bibitem{a} Referentie

\end{thebibliography}
\end{document} 